\documentclass[11pt]{article}
\usepackage{fontspec}
\usepackage[french,english]{babel}
\usepackage{amsmath}
\usepackage{amssymb}
\usepackage{amsthm}
\usepackage{mathrsfs}
\usepackage{tikz-cd}
\usepackage{enumitem}

\newcommand{\printedpage}[1]{\marginpar{\footnotesize p.~#1}}
\newcommand{\msp}[1]{\marginpar{\footnotesize [#1]}}
\newcommand{\Hot}{\mathrm{Hot}}
\newcommand{\Cat}{\mathrm{Cat}}
\newcommand{\Ob}{\operatorname{Ob}}

\begin{document}

% ===================== page 1 =====================
\printedpage{--}
% head: (none) -- title page, no page number printed

\begin{center}
\vspace*{5cm}
{\LARGE\textsf{THE DERIVATORS}}

\vspace{1.5em}
{\large\textbf{ALEXANDRE GROTHENDIECK}}

\vspace{3em}
{\large\textbf{Chapter XI}}

\vspace{0.5em}
{\large $\Hot$\textbf{-fibrations, proper functors and smooth functors etc.}\\
\textbf{(in} $\Cat$\textbf{)}}
\end{center}

\clearpage

% ===================== page 2 =====================
\printedpage{--}
% head: (none) -- no page number printed

\vspace*{6cm}
\noindent This text was deciphered and transcribed in \LaTeXe\ by M. Künzer. It was edited by M. Künzer, J. Malgoire, and G. Maltsiniotis. The transcription is as faithful as possible to the manuscript. For the few obvious corrections, or rare comments by the editors, as well as for the original numbering of the manuscript pages, typewriter characters \texttt{[typewriter]} in brackets are used. A question mark in brackets means that one is not sure of the preceding word.

\noindent This edition is provisional. Remarks, comments and corrections are welcome. Send a message to:

\noindent\textbf{[e-mail address removed]}

\hfill G. Maltsiniotis\hspace*{2.5em}

\clearpage

% ===================== page 3 =====================
\printedpage{1}
% head: A. Grothendieck : Dérivateurs, Ch. XI, Édition des Universités de Montpellier II et Paris VII
% foot: Version 28 mai 2002

\noindent\texttt{[There are no pages 1-126]}

\msp{page 127}

\section*{1\quad $\Hot$-fibrations in $\Cat$}

\noindent\textbf{Theorem 1.} \textit{Let $f : X \longrightarrow Y$ be in} $\Cat$\textit{. Equivalent conditions:}

\begin{itemize}[leftmargin=3em]
\item[\textit{a)}] \textit{For every $Y'' \xrightarrow{p} Y'$ over $Y$, whence $X'' = X \times_Y Y'' \xrightarrow{q} X' = X \times_Y Y'$, if $p$ is a} $\Hot$\textit{-equivalence, so is $q$.}

\item[\textit{a$'$)}] \textit{As in a), with $Y'' = e$, $Y'$ having a final object.}

\item[\textit{a$''$)}] \textit{As in a), with $Y'' = e$, $Y'$ having an initial object} (\footnote{keep a$''$) and b$''$), conditions dual to a$'$), b$'$) \texttt{[a$''$) and b$''$) were crossed out]}.})\textit{.}

\item[\textit{b$'$)}] \textit{For every $y \in \Ob Y$, $X_y \longrightarrow X/y$ is a} $\Hot$\textit{-equivalence, for every $y \longrightarrow y'$ in $Y$, $X/y \longrightarrow X/y'$ is a} $\Hot$\textit{-equivalence; finally these conditions still hold for $f' : X' \longrightarrow Y'$, for every base change $Y' \longrightarrow Y$ in} $\Cat$\textit{. (In fact, it suffices to require this for the first of the two conditions, concerning $X_y \longrightarrow X/y$.)}

\item[\textit{b$''$)}] \textit{Dual of b$'$), with $X_y \longrightarrow y\backslash X$ and $y'\backslash X \longrightarrow y\backslash X$.}

\item[\textit{c)}] \textit{For every locally constant sheaf of sets $F$ on $X$, $f_*(F)$ is locally constant and is computed ``fiber by fiber'', likewise for $R^1 f_*(F)$ if $F$ is a sheaf of groups, and for the $R^i f_*(F)$ if $F$ is an abelian sheaf. Moreover, we want these conditions to remain satisfied after every base change.} (\footnote{\textit{c$'$) (homological version). 1$^o$) in terms of direct images of locally constant groupoids, and 2$^o$) in terms of the $R^i f_!(L)$} \texttt{[rather $L_i f_!(L)$]}\textit{, for $L$ a locally constant abelian sheaf.}})
\end{itemize}

\noindent\textsc{Proof}. We have the $\pm$ tautological implications
\[
\mathrm{a} \implies \mathrm{a}' \implies \mathrm{b}' \implies \mathrm{c},
\]
and it remains to prove that c $\implies$ a. Since a is self-dual, being equivalent to a$'$, b$'$, it is also equivalent to the conditions dual to these, a$''$ and b$''$ (cf. cor. 2, page 75). We are reduced to proving that if c) is satisfied, and if we have
\[
\begin{tikzcd}[row sep=large, column sep=large]
X \arrow[d, "f"'] & X' \arrow[l, "q"'] \arrow[d, "f'"] \\
Y & Y' \arrow[l, "p"]
\end{tikzcd}
\]
\printedpage{2}
% head: A. Grothendieck : Dérivateurs, Ch. XI, Édition des Universités de Montpellier II et Paris VII
% foot: Version 28 mai 2002
cartesian in $\Cat$, if $p$ is a $\Hot$-equivalence, so is $q$. We use the Artin-Mazur criterion.

% ===================== page 4 (printed 2) =====================

\noindent 1$^o$) If $F$ is a locally constant sheaf of sets on $X$, we must prove that $H^0(X,F) \xrightarrow{\sim} H^0(X', q^*F)$.

\msp{page 128}

\noindent But
\[
H^0(X,F) \simeq H^0(Y,\ \underbrace{f_*(F)}_{\substack{\text{locally constant}\\ \text{by hypothesis c)}}}\ )
\qquad
\underbrace{\xrightarrow{\qquad\sim\qquad}}_{\substack{\text{since } p\\ \text{Hot-equivalence}}}
\qquad
H^0(Y', p^* f_*(F)),
\]
but
\[
p^* f_*(F) \xrightarrow{\ \sim\ } f'_*(q^*(F)),
\]
since $f_*(F)$ and $f'_*(q^*F)$ are computed ``fiber by fiber'' by hypothesis. Thus
\[
H^0(X,F) \simeq H^0(Y', f'_*(q^*(F))) \simeq H^0(X', q^*(F)),
\]

\hfill q.e.d.

\noindent 2$^o$) If $F$ is a locally constant sheaf of groups on $X$, we must prove % [sic] faiscau -> faisceau, misspelling
\[
H^1(X,F) \xrightarrow{\ \sim\ } H^1(X', q^*F).
\]
The fact that it is injective already follows from 1$^o$), which implies that the functor
\[
\begin{array}{rcl}
F\text{-torsors on } X & \longrightarrow & q^*(F)\text{-torsors on } X' \\
P & \longmapsto & f^*P % [sic] f^*P -> q^*P, the functor is the inverse image by q from X to X'
\end{array}
\]
is fully faithful, by the well-known argument using the $\mathcal{H}$om sheaves. We must see that this functor is essentially surjective.

\msp{page 129}

\noindent We have a homomorphism of exact sequences of pointed sets
\[
\begin{tikzcd}[column sep=large, row sep=large,
  every matrix/.append style={name=m},
  execute at end picture={
    \coordinate (lb) at ([xshift=0.4em]m-2-3.east);
    \coordinate (lt) at ([xshift=1.8em]m-2-3.east);
    \draw (lt |- m-1-3.north) -- (lb |- m-1-3.north) -- (lb |- m-2-3.south) -- (lt |- m-2-3.south);
    \coordinate (rb) at ([xshift=0.6em]m-2-4.east);
    \coordinate (rt) at ([xshift=-0.6em]m-2-4.east);
    \draw (rt |- m-1-3.north) -- (rb |- m-1-3.north) -- (rb |- m-2-3.south) -- (rt |- m-2-3.south);
  }]
H^1(Y, f_*F) \arrow[r, "u"] \arrow[d, "\alpha"] & H^1(X,F) \arrow[r, "v"] \arrow[d, "\beta"] & \overbrace{H^0(Y, R^1 f_*(F))}^{\ni\, \eta} \arrow[r, "w"] \arrow[d, "\gamma"] & H^2(Y, f_*F) \arrow[d, "\delta"] \\
H^1(Y', f'_*F) \arrow[r, "u'"] & \underbrace{H^1(X', F')}_{\ni\, \xi'} \arrow[r, "v'"] & \underbrace{H^0(Y', R^1 f'_*(F'))}_{\ni\, \eta'} \arrow[r, "w'"] & H^2(Y', f'_*F') % [sic] H^1(Y', f'_*F) -> H^1(Y', f'_*F'), the other entries of the lower row are in F'
\end{tikzcd}
\]
(\footnote{where we have set $F' = q^*(F)$.}), where the sheaves $f_*F$, $R^1 f_*(F)$ are locally constant by hypothesis and are computed fiber by fiber, the sheaves $f'_*(F')$ and $R^1 f'_*(F')$ also, which implies that the latter are
\printedpage{3}
% head: A. Grothendieck : Dérivateurs, Ch. XI, Édition des Universités de Montpellier II et Paris VII
% foot: Version 28 mai 2002
canonically isomorphic (by the base change arrows) to the inverse images of their counterparts $f_*(F)$, $R^1 f_*(F)$. Since $p$ is a $\Hot$-equivalence, this implies that the arrows $\alpha$, $\gamma$ are \texttt{[some]} isomorphisms. Thus if $\xi' \in H^1(X', F')$, its image $\eta'$ in $H^0(Y', R^1 f'_*(F'))$ comes from a well-determined $\eta \in H^0(Y, R^1 f_*(F))$. We must see that this $\eta$ comes from a $\xi \in H^1(X,F)$. This (if I remember correctly) is possible if and only if a certain obstruction $w(\eta)$ in $H^2(Y, f_*(F))$ is in the neutral subset of $H^2$. \textbf{NB} The stack on $Y^\wedge$ of $F$-torsors on the $X^\wedge/U$ that lift $\eta$ is an $\mathcal{F} = f_*F$-gerbe $\mathcal{G}_\eta$, and we must see that it is neutral. But there is something better to do:

% ===================== page 5 (printed 3) =====================

\msp{page 130}

\noindent We note that this gerbe $\mathcal{G}$ is locally constant (since it corresponds to a locally constant band $f_*(G)$) and that $\mathcal{G}'$ is its inverse image. But since $p$ is a $\Hot$-equivalence, it follows that the functor
\[
\underline{\Gamma}(Y, \mathcal{G}) \longrightarrow \underline{\Gamma}(Y', \mathcal{G}')
\]
is an equivalence of categories, so $\xi'$ comes (up to isomorphism) from a unique $\xi$ of $\underline{\Gamma}(Y, \mathcal{G})$, q.e.d.

\noindent 3$^o$) Let $F$ be a locally constant abelian sheaf on $X$, we must see \texttt{[that]}
\[
H^i(X,F) \xrightarrow{\ \sim\ } H^i(X', F') \qquad (F' = q^*F), \tag*{$(*)$}
\]
for every $i \in \mathbf{N}$. But the two terms are the abutment of the Leray spectral sequences, with initial terms
\[
\begin{aligned}
E_2^{pq} &\simeq H^p(Y, R^q f_*(F)) \\
E^{\prime pq}_2 &\simeq H^p(Y', R^q f'_*(F')),
\end{aligned}
\]
and $(*)$ is associated with a homomorphism of spectral sequences, which for the initial terms is
\[
H^p(Y, R^q f_*(F)) \xrightarrow{\sim} H^p(Y', R^q f'_*(F)), \tag*{$(**)$} % [sic] R^q f'_*(F) -> R^q f'_*(F'), cf. the E'_2 term just above
\]
associated with
\[
p^*(R^q f_*(F)) \longrightarrow R^q f'_*(F'). \tag*{$(*{*}*)$}
\]
But by hypothesis the $R^q f_*(F)$, $R^q f'_*(F')$ are locally constant and are computed fiber by fiber,

\msp{page 131}

\noindent whence it follows that the base change homomorphisms $(*{*}*)$ are \texttt{[some]} isomorphisms, and since we are dealing with locally constant sheaves and $p$ is a $\Hot$-equivalence, it follows that $(**)$ is \texttt{[an]} isomorphism, hence also $(*)$ on the abutments, q.e.d.

\noindent\textbf{Definition 1.} Under the equivalent conditions of theorem 1, we say that $f$ is a $\Hot$-\textit{fibration.}

\end{document}
